\section*{Capítulo \ref{ch:b3:banach} --- Espaços de Banach e os
teoremas fundamentais}

\begin{solution}{exo:b3:banach:1}
(a) $\norm{Sx}_2 = \norm x_2$: $S$ é uma isometria, $\vertiii S =
1$. Para o deslocamento à esquerda: $\norm{S^*x}_2^2 = \sum_{n \geq
2}\abs{x_n}^2 \leq \norm x_2^2$, com igualdade para $x = e_2$:
$\vertiii{S^*} = 1$.

(b) $\norm{M_ax}_2^2 = \sum\abs{a_n}^2\abs{x_n}^2 \leq \norm
a_\infty^2\norm x_2^2$; testar $x = e_n$ dá $\vertiii{M_a}
\geq \abs{a_n}$ para todo $n$: $\vertiii{M_a} = \norm a_\infty$.

(c) $\abs{\Lambda(f)} \leq \int_0^1\abs f \leq \norm f_\infty$:
$\norm\Lambda \leq 1$. Para $\varepsilon > 0$, seja $f_\varepsilon$
igual a $1$ em $[0, \frac12 - \varepsilon]$, $-1$ em $[\frac12 +
\varepsilon, 1]$, afim entre eles: $\norm{f_\varepsilon}_\infty =
1$ e $\Lambda(f_\varepsilon) \geq 1 - 2\varepsilon$:
$\norm\Lambda = 1$. Não atingida: $\Lambda(f) = 1$ com $\norm
f_\infty \leq 1$ força $\int_0^{1/2}f = \frac12$ e
$\int_{1/2}^1 f = -\frac12$, isto é (\omterm{def:b3:topology:continuity}{continuidade}, $\abs f \leq 1$),
$f \equiv 1$ em $[0, \frac12]$ e $f \equiv -1$ em $[\frac12,
1]$: contradição em $\frac12$.
\end{solution}

\begin{solution}{exo:b3:banach:2}
Escreva $S = T\bigl(I - T^{-1}(T - S)\bigr)$ com
$\vertiii{T^{-1}(T-S)} \leq \vertiii{T^{-1}}\,\vertiii{T - S} =
\theta < 1$: pela~\cref{prop:b3:banach:neumann}, $S$ é inversível
com $S^{-1} = \sum_{n\geq0}\bigl(T^{-1}(T -
S)\bigr)^nT^{-1}$, donde
\[
\vertiii{S^{-1} - T^{-1}} \leq
\sum_{n\geq1}\theta^n\,\vertiii{T^{-1}}
= \frac{\vertiii{T^{-1}}^2\,\vertiii{T - S}}{1 - \theta}.
\]
Para o sistema linear: $x_T = T^{-1}b$ e $x_S = S^{-1}b$
diferem, no máximo, dessa cota vezes $\norm b$ --- pequenas
perturbações de um sistema inversível permanecem unicamente \omterm{def:b3:groups:derived}{solúveis},
com dependência lipschitziana da solução em relação ao operador.
\end{solution}

\begin{solution}{exo:b3:banach:3}
(a) $\ell^1$: seja $(x^{(k)})$ de Cauchy. Cada coordenada é
de Cauchy ($\abs{x^{(k)}_n - x^{(l)}_n} \leq \norm{x^{(k)} -
x^{(l)}}_1$): seja $x_n = \lim_kx^{(k)}_n$. Dado $\varepsilon$,
para $k, l \geq K$: $\sum_{n \leq N}\abs{x^{(k)}_n - x^{(l)}_n}
\leq \varepsilon$ para todo $N$; faça $l \to \infty$ e depois $N \to
\infty$: $\norm{x^{(k)} - x}_1 \leq \varepsilon$, e $x =
x^{(k)} - (x^{(k)} - x) \in \ell^1$. $\ell^\infty$: ser de Cauchy para
$\norm\cdot_\infty$ é ser uniformemente de Cauchy: converge uniformemente para
uma sequência limitada. $c_0$ é fechado em $\ell^\infty$: se
$x^{(k)} \to x$ uniformemente com $x^{(k)}_n \to_n 0$, então
$\abs{x_n} \leq \norm{x - x^{(k)}}_\infty + \abs{x^{(k)}_n}$
dá $\limsup_n\abs{x_n} \leq \varepsilon$: $x \in c_0$; um
subespaço fechado de um espaço de Banach é de Banach. Sequências finitas:
seu \omterm{def:b3:topology:interior}{fecho} contém todo $x \in c_0$ (as truncaturas convergem:
$\sup_{n>N}\abs{x_n} \to 0$) e está contido no fechado
$c_0$.

(b) Por homogeneidade, suponha $\norm x_p = 1$: então $\abs{x_n} \leq
1$ para todo $n$, de modo que $\abs{x_n}^q \leq \abs{x_n}^p$ e $\norm
x_q \leq 1 = \norm x_p$; para $q = \infty$, $\abs{x_n} \leq
\norm x_p$ diretamente. Estrita: $x_n = n^{-\alpha}$ com
$\frac1q < \alpha \leq \frac1p$ está em $\ell^q \setminus
\ell^p$ (série de Riemann).
\end{solution}

\begin{solution}{exo:b3:banach:4}
($\leq$) Se $\norm f \leq 1$ e $f\restriction_F = 0$: para
todo $y \in F$, $\abs{f(x)} = \abs{f(x - y)} \leq \norm{x -
y}$; tome o ínfimo. ($\geq$, atingido)
O~\cref{cor:b3:banach:hbnormed}(3) produz $f$ com
$f\restriction_F = 0$, $\norm f \leq 1$, $f(x) = d(x, F)$: o
supremo é um máximo. Consequência: $F \subseteq \bigcap\{\ker
f : f\restriction_F = 0\}$ trivialmente, e um ponto $x \notin F$
é excluído da interseção pelo funcional acima
($f(x) = d(x,F) > 0$, sendo $F$ fechado).
\end{solution}

\begin{solution}{exo:b3:banach:5}
Para $y \in \ell^\infty$: $\abs{\Lambda_y(x)} =
\abs{\sum x_ny_n} \leq \norm y_\infty\norm x_1$, de modo que
$\norm{\Lambda_y} \leq \norm y_\infty$; testando em $e_n$:
$\abs{y_n} = \abs{\Lambda_y(e_n)} \leq \norm{\Lambda_y}$:
igualdade. Reciprocamente, dado $\Lambda \in (\ell^1)'$, ponha $y_n =
\Lambda(e_n)$: $\abs{y_n} \leq \norm\Lambda$, de modo que $y \in
\ell^\infty$; $\Lambda$ e $\Lambda_y$ coincidem nas sequências
finitas, densas em $\ell^1$ ($\norm{x - \sum_{n\leq
N}x_ne_n}_1 = \sum_{n>N}\abs{x_n} \to 0$): $\Lambda =
\Lambda_y$. A aplicação $y \mapsto \Lambda_y$ é linear, isométrica e
sobrejetora. Para $(\ell^\infty)'$, o mesmo início produz uma sequência
$y_n = \Lambda(e_n)$, mas as sequências finitas \emph{não} são densas
em $\ell^\infty$
(a sequência constante $\mathbf 1$ está à distância $1$ de todas
elas), de modo que $\Lambda$ não fica determinada pelos $y_n$ --- e,
de fato, $(\ell^\infty)' \neq \ell^1$
(\cref{pb:b3:banach:1}).
\end{solution}

\begin{solution}{exo:b3:banach:6}
Para cada $x$ fixo, $y \mapsto B(x, y)$ é linear \omterm{def:b3:topology:continuity}{contínua}:
$\sup_{\norm y \leq 1}\norm{B(x,y)} < \infty$. Logo a família
$\{B(\cdot, y) : \norm y \leq 1\} \subseteq \mathcal L(E, G)$
(cada membro \omterm{def:b3:topology:continuity}{contínuo}, pela \omterm{def:b3:topology:continuity}{continuidade} em $x$) é pontualmente
limitada no espaço de Banach $E$: Banach--Steinhaus
(\cref{thm:b3:banach:banachsteinhaus}) fornece $C$ com
$\norm{B(x,y)} \leq C\norm x$ para todo $\norm y \leq 1$;
a homogeneidade em $y$ conclui: $\norm{B(x,y)} \leq
C\norm x\norm y$.
\end{solution}

\begin{solution}{exo:b3:banach:7}
(a) $\norm f_1 \leq \norm f_\infty$: a identidade é limitada
e bijetora. Sua inversa é ilimitada: $f_n(x) = x^n$ tem
$\norm{f_n}_1 = \frac1{n+1} \to 0$, mas $\norm{f_n}_\infty = 1$.
Nenhuma contradição com o~\cref{cor:b3:banach:equivnorms}: $(\mathcal
C(\intcc01), \norm\cdot_1)$ \emph{não} é \omterm{def:b3:complete:complete}{completo}
(\cref{exo:b3:complete:1}); o corolário exige \omterm{def:b3:complete:complete}{completude}
dos dois lados.

(b) No espaço $E_0$ das sequências finitas (denso em $\ell^2$),
$\varphi(x) = \sum_n n\,x_n$ é linear e ilimitado
($\varphi(e_n) = n$ com $\norm{e_n}_2 = 1$). O teorema do gráfico
fechado não se aplica: $E_0$ não é \omterm{def:b3:complete:complete}{completo} --- e $\varphi$
não tem extensão \omterm{def:b3:topology:continuity}{contínua} a $\ell^2$, ilustrando que
densidade sem \omterm{def:b3:topology:continuity}{continuidade} uniforme nada pode
(\cref{thm:b3:complete:extension}).
\end{solution}

\begin{solution}{exo:b3:banach:8}
Verificamos a hipótese do gráfico fechado. Sejam $x_k \to x$ e
$Tx_k \to z$ em $\ell^2$. Para todo $y$:
\[
\langle z, y\rangle = \lim_k\langle Tx_k, y\rangle
= \lim_k \langle x_k, Ty\rangle = \langle x, Ty\rangle
= \langle Tx, y\rangle,
\]
usando a \omterm{def:b3:topology:continuity}{continuidade} do produto interno em cada entrada
(Cauchy--Schwarz) e a simetria duas vezes. Logo $z - Tx$ é ortogonal
a todo $y$, em particular a si mesmo: $z = Tx$. O gráfico é
fechado e $\ell^2$ é de Banach: $T$ é limitado
(\cref{thm:b3:banach:closedgraph}). Assim, um operador simétrico
definido em \emph{todo} $\ell^2$ é automaticamente limitado;
operadores simétricos genuinamente ilimitados (posição, momento,
hamiltonianos) devem viver em subespaços densos próprios.
\end{solution}

\begin{solution}{exo:b3:banach:9}
Primeiro, $\norm{\Lambda_n} = \sum_i\abs{w_{i,n}}$: $\leq$ é a
desigualdade triangular; $\geq$ testando com uma $f$ linear por partes
com $\norm f_\infty \leq 1$ e $f(x_{i,n}) =
\operatorname{sign}(w_{i,n})$ (interpole linearmente entre os
finitos nós; onde os nós coincidem, os sinais concordam).

($\Rightarrow$) A convergência pontual em toda $f$ implica (i)
e a limitação pontual, de modo que Banach--Steinhaus
(\cref{thm:b3:banach:banachsteinhaus}) no espaço de Banach $\mathcal
C(\intcc01)$ dá (ii).

($\Leftarrow$) Seja $M = \sup_n\norm{\Lambda_n} + 1$. Dados $f$
e $\varepsilon$, escolha um polinômio $P$ com $\norm{f -
P}_\infty < \varepsilon/(2M)$
(\cref{cor:b3:complete:weierstrass}); então
\[
\Bigl|\Lambda_n f - \int_0^1 f\Bigr|
\leq \abs{\Lambda_n(f - P)} + \Bigl|\Lambda_nP -
\int_0^1P\Bigr| + \Bigl|\int_0^1(P - f)\Bigr|
\leq \varepsilon + \Bigl|\Lambda_nP - \int_0^1P\Bigr| \to
\varepsilon .
\]
Pesos positivos, exatidão nas constantes: $\sum_i\abs{w_{i,n}}
= \sum_iw_{i,n} = \Lambda_n(\mathbf 1) = \int_0^1 1 = 1$ para
regras exatas nas constantes --- (ii) vale com constante $1$.
\end{solution}

\begin{solution}{exo:b3:banach:10}
A linearidade de $J$ é formal; $\norm{J(x)} = \sup_{\norm f \leq
1}\abs{f(x)} = \norm x$ pelo~\cref{cor:b3:banach:hbnormed}(2). Se $\dim E = n$: $\dim E' =
n$ (uma base dá os funcionais coordenados), de modo que $\dim E'' = n$,
e o $J$ injetor (isométrico) é sobrejetor. Separabilidade: seja
$(f_n)$ denso em $E'$ e escolha $\norm{x_n} = 1$ com
$\abs{f_n(x_n)} \geq \frac12\norm{f_n}$. Seja $F =
\overline{\operatorname{Vect}}(x_n)$; se $F \neq E$, tome $g
\in E'$, $g \neq 0$, nulo em $F$
(\cref{cor:b3:banach:hbnormed}(3)); escolha $f_{n_k} \to g$:
\[
\norm{f_{n_k} - g} \geq \abs{(f_{n_k} - g)(x_{n_k})}
= \abs{f_{n_k}(x_{n_k})} \geq \tfrac12\norm{f_{n_k}}
\geq \tfrac12\bigl(\norm g - \norm{f_{n_k} - g}\bigr),
\]
de modo que $\norm{f_{n_k} - g} \geq \frac13\norm g > 0$: contradição.
Logo $F = E$, e as combinações racionais (ou $\Q + \iu\Q$) dos
$x_n$ formam um conjunto enumerável denso.
\end{solution}

\begin{solution}{exo:b3:banach:11}
(a) Boa definição: $d(x, F)$ depende apenas de $\bar x$
(transladar $x$ por $F$ não muda a distância).
A homogeneidade e a desigualdade triangular passam de $\norm\cdot$
pelo ínfimo. A separação precisa que o subespaço seja fechado: $\norm{\bar
x} = 0$ significa $d(x, F) = 0$, isto é, $x \in \bar F = F$, isto é,
$\bar x = 0$. $\vertiii\pi \leq 1$: $\norm{\bar x} \leq
\norm x$. Bola aberta sobre bola aberta: se $\norm{\bar x} < 1$,
algum representante tem $\norm{x - y} < 1$; reciprocamente,
$\pi(B_E(0,1)) \subseteq B_{E/F}(0,1)$ pela desigualdade entre
normas --- de modo que $\pi$ é aberta, o caso modelo do teorema da
aplicação aberta.

(b) Seja $\sum_k\norm{\bar x_k} < \infty$ e levante com
$\norm{x_k} \leq \norm{\bar x_k} + 2^{-k}$: então
$\sum\norm{x_k} < \infty$, de modo que $s = \sum_kx_k$ converge no
espaço de Banach $E$ (\cref{exo:b3:complete:1}(b)), e a \omterm{def:b3:topology:continuity}{continuidade}
de $\pi$ dá $\sum_k\bar x_k = \bar s$: toda série absolutamente
convergente de $E/F$ converge, o que é equivalente à
\omterm{def:b3:complete:complete}{completude} (mesmo exercício).

(c) A aplicação $\lambda(x) = \lim_nx_n$ é linear $c \to K$,
anula-se exatamente em $c_0$ e, portanto, induz uma bijeção linear
$c/c_0 \to K$. Isometria: $d(x, c_0) = \abs{\lambda(x)}$ ---
$\leq$: subtraia de $x$ a sequência $x - \lambda(x)\mathbf
1 \in c_0$, deixando $\lambda(x)\mathbf 1$ de norma
$\abs{\lambda(x)}$; $\geq$: para $y \in c_0$, $\norm{x -
y}_\infty \geq \limsup_n\abs{x_n - y_n} = \abs{\lambda(x)}$.
\end{solution}

\begin{solution}{exo:b3:banach:12}
(a) $W = \ker P$ é fechado (pré-imagem de $0$ por uma
\omterm{def:b3:topology:continuity}{aplicação contínua}); $V = \operatorname{im}P = \ker(I - P)$
(de fato, $Px = x$ se, e somente se, $x \in \operatorname{im}P$, usando $P^2 =
P$), fechado do mesmo modo. Todo $x$ se reparte como $Px + (x - Px)$
com $Px \in V$, $x - Px \in W$, e $V \cap W = 0$ ($x = Px
= 0$): $E = V \oplus W$.

(b) O argumento do gráfico: sejam $x_n \to x$ e $Px_n \to z$.
Então $z \in V$ ($V$ fechado, $Px_n \in V$) e $x_n - Px_n \to
x - z \in W$ ($W$ fechado). Logo $x = z + (x - z)$ com $z \in
V$, $x - z \in W$; pela unicidade da decomposição, $z =
Px$. O gráfico de $P$ é fechado e $E$ é de Banach: $P$ é
limitado (\cref{thm:b3:banach:closedgraph}).

(c) (a) e (b) juntos são a equivalência. Num espaço de Hilbert,
todo $V$ fechado tem o complemento fechado $V^\perp$
(\cref{ch:b3:hilbert}): todo subespaço fechado é
complementado. Em espaços de Banach gerais isso falha --- $c_0$
não tem complemento fechado em $\ell^\infty$ (teorema de
Phillips, além de nossas ferramentas) ---, de modo que projeções limitadas são um
privilégio, e o teorema do gráfico fechado é exatamente o que
converte decomposições geométricas em operadores limitados.
\end{solution}

\begin{solution}{pb:b3:banach:1}
\textbf{1.} Para $a, b > 0$: pela concavidade de $\log$,
$\log\bigl(\tfrac{a^p}p + \tfrac{b^q}q\bigr) \geq \tfrac1p\log
a^p + \tfrac1q\log b^q = \log(ab)$; exponencie. (Se $ab = 0$,
a desigualdade é trivial.) Igualdade se, e somente se, $a^p = b^q$.

\textbf{2.} Podemos supor $\norm x_p = \norm y_q = 1$
(homogeneidade; os casos nulos são triviais). Então
\[
\abs{\langle x, y\rangle} \leq \sum_n\abs{x_n}\abs{y_n}
\leq \sum_n\Bigl(\frac{\abs{x_n}^p}p +
\frac{\abs{y_n}^q}q\Bigr) = \frac1p + \frac1q = 1 =
\norm x_p\norm y_q .
\]
A igualdade exige $\abs{x_n}^p = \abs{y_n}^q$ para todo $n$
(caso de igualdade em Young) e o alinhamento das fases de
$x_ny_n$.

\textbf{3.} $\abs{x_n + y_n}^p \leq \bigl(\abs{x_n} +
\abs{y_n}\bigr)\abs{x_n + y_n}^{p-1}$; somando e aplicando
Hölder ($p$ contra $q$, notando $(p - 1)q = p$):
\[
\norm{x + y}_p^p \leq \bigl(\norm x_p + \norm
y_p\bigr)\,\Bigl(\sum_n\abs{x_n + y_n}^{p}\Bigr)^{1/q}
= \bigl(\norm x_p + \norm y_p\bigr)\,\norm{x+y}_p^{p/q};
\]
se $\norm{x + y}_p \neq 0$ (senão é trivial), divida por
$\norm{x+y}_p^{p/q}$ e use $p - \frac pq = 1$. (Finitude de
$\norm{x+y}_p$ antes: $\abs{x_n+y_n}^p \leq
2^p(\abs{x_n}^p + \abs{y_n}^p)$.)

\textbf{4.} \omterm{def:b3:complete:complete}{Completude}: como para $\ell^1$
(\cref{exo:b3:banach:3}), limites coordenada a coordenada mais a
cota uniforme das caudas $\sum_{n\leq N}\abs{x^{(k)}_n - x^{(l)}_n}^p
\leq \varepsilon^p$, fazendo $l$ e depois $N$ tenderem ao infinito.
Densidade das sequências finitas: $\norm{x - \sum_{n\leq
N}x_ne_n}_p^p = \sum_{n > N}\abs{x_n}^p \to 0$.

\textbf{5.} Hölder dá $\abs{\Lambda_y(x)} \leq \norm
x_p\norm y_q$: $\norm{\Lambda_y} \leq \norm y_q$. Testando: seja
$x^{(N)}_n = \abs{y_n}^{q-1}\overline{\operatorname{sign}}(y_n)$
para $n \leq N$, $0$ além disso (com $\operatorname{sign}$ a
fase unimodular, de modo que $x_ny_n = \abs{y_n}^q$). Então
$\Lambda_y(x^{(N)}) = \sum_{n\leq N}\abs{y_n}^q$ e
$\norm{x^{(N)}}_p = \bigl(\sum_{n\leq
N}\abs{y_n}^{q}\bigr)^{1/p}$ (pois $(q-1)p = q$), de modo que
\[
\norm{\Lambda_y} \geq \Bigl(\sum_{n\leq
N}\abs{y_n}^q\Bigr)^{1 - 1/p} \xrightarrow[N\to\infty]{}
\norm y_q .
\]

\textbf{6.} Ponha $y_n = \Lambda(e_n)$. Com os mesmos vetores de
teste, $\sum_{n \leq N}\abs{y_n}^q = \Lambda(x^{(N)}) \leq
\norm\Lambda\,\bigl(\sum_{n\leq N}\abs{y_n}^q\bigr)^{1/p}$,
donde $\bigl(\sum_{n\leq N}\abs{y_n}^q\bigr)^{1/q} \leq
\norm\Lambda$ para todo $N$: $y \in \ell^q$, $\norm y_q \leq
\norm\Lambda$. Os funcionais $\Lambda$ e $\Lambda_y$ coincidem
nas sequências finitas densas (questão 4): $\Lambda =
\Lambda_y$. Com a questão 5, $y \mapsto \Lambda_y$ é um
isomorfismo isométrico $\ell^q \cong (\ell^p)'$.

\textbf{7.} Seja $\xi \in (\ell^p)''$. Compondo com a
isometria $\ell^q \cong (\ell^p)'$ da questão 6, $\xi$ define
um elemento de $(\ell^q)'$, que (questão 6 com $p, q$
trocados) é $\Lambda_z$ para um único $z \in \ell^p$: para todo
$y \in \ell^q$, $\xi(\Lambda_y) = \sum_nz_ny_n$. Por outro
lado, $J(z)(\Lambda_y) = \Lambda_y(z) = \sum_ny_nz_n$: o mesmo
valor. Como todo elemento de $(\ell^p)'$ é algum $\Lambda_y$,
$\xi = J(z)$: $J$ é sobrejetora --- $\ell^p$ é \omterm{rem:b3:banach:bidual}{reflexivo}.

\textbf{8.} As sequências finitas com entradas racionais (parte real e
imaginária) são enumeráveis e densas em $\ell^p$ ($p < \infty$) e
em $c_0$. Em $\ell^\infty$: a família $\{\mathbf 1_A : A
\subseteq \N\}$ é não enumerável, com $\norm{\mathbf 1_A -
\mathbf 1_B}_\infty = 1$ para $A \neq B$; as bolas $B(\mathbf
1_A, \frac12)$ são duas a duas disjuntas, e um \omterm{def:b3:topology:interior}{conjunto denso} deve encontrar
cada uma: não existe \omterm{def:b3:topology:interior}{conjunto denso} enumerável.

\textbf{9.} Se $\ell^1$ fosse \omterm{rem:b3:banach:bidual}{reflexivo}, então $(\ell^\infty)'
\cong ((\ell^1)')' = J(\ell^1)$ seria separável (imagem isométrica
do separável $\ell^1$); pelo~\cref{exo:b3:banach:10}, a separabilidade do \omterm{def:b3:banach:operator}{dual}
$(\ell^\infty)'$ forçaria $\ell^\infty$ separável ---
contradizendo a questão 8. Logo $\ell^1$ não é \omterm{rem:b3:banach:bidual}{reflexivo} (e
$(\ell^\infty)'$ é estritamente maior que $\ell^1$, como a questão
11 torna concreto).

\textbf{10.} A homogeneidade de $p$ é clara; a subaditividade:
as médias são lineares, e $\limsup(u_n + v_n) \leq \limsup u_n +
\limsup v_n$. Em $c$: as médias de Cesàro de uma convergente
convergem para seu limite, de modo que $p(x) = \lim x =
\mathrm{LIM}(x)$ ali; em particular, $\mathrm{LIM} \leq p$ em
$c$. Hahn--Banach (\cref{thm:b3:banach:hahnbanach}) estende
$\mathrm{LIM}$ a $\ell^\infty_\R$ com $\mathrm{LIM} \leq p$
globalmente. Positividade: para $x \geq 0$, $-\mathrm{LIM}(x) =
\mathrm{LIM}(-x) \leq p(-x) = \limsup\text{méd}(-x) \leq 0$.
Invariância por deslocamento: as médias de $x - Sx$ telescopam em
$\frac{x_1 - x_{n+1}}n \to 0$, de modo que $p(\pm(x - Sx)) = 0$ e
$\mathrm{LIM}(x - Sx) = 0$. Cotas: $\mathrm{LIM}(x) \leq p(x)
\leq \limsup x$ (as médias ficam atrás dos supremos), e aplicar isso
a $-x$ dá a cota inferior.

\textbf{11.} $e_n \in c_0$, de modo que $\mathrm{LIM}(e_n) = 0$ para todo
$n$. Se $\mathrm{LIM} = \Lambda_y$ com $y \in \ell^1$, então
$y_n = \Lambda_y(e_n) = 0$ para todo $n$: $\Lambda_y = 0$; mas
$\mathrm{LIM}(\mathbf 1) = 1$. Logo $\mathrm{LIM} \in
(\ell^\infty)' \setminus \{\Lambda_y : y \in \ell^1\}$:
$(\ell^\infty)' \neq \ell^1$, de novo.

\textbf{12.} Para $x = (0,1,0,1,\dots)$: $x + Sx = \mathbf 1$,
de modo que $2\,\mathrm{LIM}(x) = \mathrm{LIM}(x) + \mathrm{LIM}(Sx) =
1$: $\mathrm{LIM}(x) = \frac12$. Se $\varphi$ fosse uma
extensão multiplicativa e invariante por deslocamento do limite: $x\cdot
Sx = 0$ dá $\varphi(x)\varphi(Sx) = \varphi(x)^2 = 0$, de modo que
$\varphi(x) = 0$; mas $\varphi(x) + \varphi(Sx) =
\varphi(\mathbf 1) = 1$ dá $2\varphi(x) = 1$: contradição.
Média e multiplicação não podem coexistir.

\textbf{13.} Seja $\norm{x^{(k)}}_p \leq M$. As coordenadas são
limitadas por $M$: uma extração diagonal dá uma subsequência
(ainda escrita $x^{(k)}$) com $x^{(k)}_n \to x_n$ para todo
$n$. Então $x \in \ell^p$: $\sum_{n\leq N}\abs{x_n}^p =
\lim_k\sum_{n\leq N}\abs{x^{(k)}_n}^p \leq M^p$ para todo $N$.
Convergência fraca: para $y \in \ell^q$ e $N$ arbitrário,
\[
\abs{\Lambda_y(x^{(k)} - x)} \leq
\Bigl|\sum_{n \leq N}(x^{(k)}_n - x_n)y_n\Bigr|
+ 2M\Bigl(\sum_{n>N}\abs{y_n}^q\Bigr)^{1/q},
\]
em que o primeiro termo tende a $0$ quando $k \to \infty$ (finitas
coordenadas) e o segundo é pequeno para $N$ grande
(Hölder na cauda): $\Lambda_y(x^{(k)}) \to \Lambda_y(x)$
para todo $y$ --- convergência fraca, já que todo funcional é um
$\Lambda_y$ (questão 6). Em $\ell^1$ isso falha: considere
$(e_n)$, limitada. Toda subsequência $(e_{n_k})$ converge
coordenada a coordenada para $0$, de modo que seu único candidato a limite fraco é $0$;
mas, testando contra $y \in \ell^\infty$ definida por $y_{n_k} =
(-1)^k$ (e $0$ nas demais), $\Lambda_y(e_{n_k}) = (-1)^k$
diverge. Nenhuma subsequência fracamente convergente: a \omterm{thm:b3:topology:metriccompact}{compacidade
sequencial} fraca das bolas caracteriza o mundo \omterm{rem:b3:banach:bidual}{reflexivo}.

\textbf{14.} Para $y \in \ell^1$, $\Lambda_y(x) = \sum x_ny_n$
está definida em $c_0$ com $\abs{\Lambda_y(x)} \leq \norm
x_\infty\norm y_1$, e testar em $x^{(N)} =
(\operatorname{sign}y_1, \dots, \operatorname{sign}y_N, 0,
\dots) \in c_0$ dá $\Lambda_y(x^{(N)}) =
\sum_{n\leq N}\abs{y_n} \to \norm y_1$:
$\norm{\Lambda_y} = \norm y_1$. Reciprocamente, para $\Lambda \in
(c_0)'$ ponha $y_n = \Lambda(e_n)$; os mesmos testes dão
$\sum_{n \leq N}\abs{y_n} = \Lambda(x^{(N)}) \leq
\norm\Lambda$, de modo que $y \in \ell^1$, e $\Lambda = \Lambda_y$ nas
sequências finitas densas, logo em toda parte. A cadeia dos
duais: $(c_0)' = \ell^1$, $(\ell^1)' = \ell^\infty$
(\cref{exo:b3:banach:10}), $(\ell^\infty)' \supsetneq \ell^1$
(questões 9--11): a reflexividade falha logo no primeiro passo
--- $c_0'' = \ell^\infty \neq c_0$ --- e nunca se recupera.

\textbf{15.} $J x^{(k)} \in E'' = (E')'$ é uma família de
funcionais limitados no espaço \emph{de Banach} $E'$ (duais
são \omterm{def:b3:complete:complete}{completos}, \cref{prop:b3:banach:llcomplete}); para cada
$\Lambda \in E'$, a sequência $Jx^{(k)}(\Lambda) =
\Lambda(x^{(k)})$ converge, logo é limitada.
Banach--Steinhaus em $E'$ dá $\sup_k\norm{Jx^{(k)}} <
\infty$, e $J$ é isométrica
(\cref{rem:b3:banach:bidual}): $\sup_k\norm{x^{(k)}} <
\infty$.

\textbf{16.} ($\Rightarrow$) A limitação é a questão 15;
as coordenadas são os funcionais $\Lambda_{e_n}$.
($\Leftarrow$) Sejam $M = \sup_k\norm{x^{(k)}}_p$, $y \in
\ell^q$, $\varepsilon > 0$; escolha $N$ com
$\bigl(\sum_{n>N}\abs{y_n}^q\bigr)^{1/q} < \varepsilon$. Então
\[
\abs{\Lambda_y(x^{(k)} - x)} \leq
\sum_{n\leq N}\abs{x^{(k)}_n - x_n}\,\abs{y_n} +
\norm{x^{(k)} - x}_p\,\varepsilon
\leq \sum_{n\leq N}\abs{x^{(k)}_n - x_n}\abs{y_n} +
(M + \norm x_p)\,\varepsilon,
\]
e a soma finita tende a $0$: $\limsup \leq (M + \norm
x_p)\varepsilon$ para todo $\varepsilon$. (Que $x \in
\ell^p$ com $\norm x_p \leq M$ decorre de cotas de seções finitas
ao estilo de Fatou: $\sum_{n \leq N}\abs{x_n}^p =
\lim_k\sum_{n\leq N}\abs{x^{(k)}_n}^p \leq M^p$.) Para $e_k$
em $\ell^2$: limitada, coordenada a coordenada $\to 0$, de modo que $e_k
\rightharpoonup 0$, mas $\norm{e_k} = 1$: a unidade de massa
escapa para índices infinitos, invisível a todo funcional
fixo.

\textbf{17.} Tome $\Lambda$ com $\norm\Lambda = 1$ e
$\Lambda(x) = \norm x$ (\cref{cor:b3:banach:hbnormed}). Então
$\norm x = \lim\Lambda(x^{(k)}) \leq
\liminf\norm\Lambda\,\norm{x^{(k)}} =
\liminf\norm{x^{(k)}}$. (Com $e_k \rightharpoonup 0$: $0
\leq \liminf 1$, e a desigualdade pode ser estrita.)

\textbf{18.} Em $\ell^2$, $\norm{x^{(k)} - x}_2^2 =
\norm{x^{(k)}}_2^2 - 2\operatorname{Re}\langle x^{(k)},
x\rangle + \norm x_2^2$ (caso real: $-2\langle x^{(k)},
x\rangle$). A convergência fraca aplicada ao funcional
$\Lambda_x$ dá $\langle x^{(k)}, x\rangle \to \norm
x_2^2$, e as normas convergem por hipótese: o membro direito
tende a $\norm x^2 - 2\norm x^2 + \norm x^2 = 0$.
Contraexemplo sem convergência das normas: $e_k
\rightharpoonup 0$, $\norm{e_k - 0} = 1 \not\to 0$.

\textbf{19.} Convergência das coordenadas: aplique os funcionais
$\Lambda_{e_n} \in (\ell^1)' = \ell^\infty$ ($y = e_n$).
Construção: tendo escolhido $k_{j-1}, N_{j-1}$, tome $k_j >
k_{j-1}$ tão grande que $\sum_{n \leq
N_{j-1}}\abs{x^{(k_j)}_n} < \frac\delta{20}$ (finitas
coordenadas, cada uma $\to 0$), e depois $N_j > N_{j-1}$ tão grande
que a cauda satisfaça $\sum_{n > N_j}\abs{x^{(k_j)}_n} <
\frac\delta{20}$ (convergência da série que define
$\norm{x^{(k_j)}}_1$). O bloco $B_j = \intoc{N_{j-1}}{N_j}$
carrega então tudo, salvo $\frac\delta{10}$, da massa de
$x^{(k_j)}$.

\textbf{20.} Com $y$ como definida ($\abs{y_n} \leq 1$
em toda parte):
\[
\langle x^{(k_j)}, y\rangle
= \sum_{n \in B_j}\abs{x^{(k_j)}_n}
+ \sum_{n \notin B_j}x^{(k_j)}_ny_n
\geq \Bigl(\norm{x^{(k_j)}}_1 - \frac\delta{10}\Bigr) -
\frac\delta{10}
\geq \delta - \frac{2\delta}{10} = \frac{4\delta}5 > 0,
\]
a desigualdade do meio porque a massa fora do bloco é no
máximo $\frac\delta{10}$ (questão 19). Mas $y \in
\ell^\infty = (\ell^1)'$ e $x^{(k)} \rightharpoonup 0$
forçam $\langle x^{(k_j)}, y\rangle \to 0$: contradição.
Logo as sequências fracamente nulas de $\ell^1$ são nulas em norma e, por
translação, as fracamente convergentes convergem em norma: o teorema de
Schur.

\textbf{21.} Uma subsequência fracamente convergente de $(e_k)$ teria
limite $0$ (coordenadas), logo, por Schur, $\norm{e_{k_j}}_1
\to 0$ --- mas as normas valem $1$. Assim, não existe subsequência fracamente
convergente, como se viu à mão na questão 13. Nenhum
paradoxo: Schur diz que as sequências não conseguem distinguir a \omterm{def:b3:topology:topology}{topologia}
fraca da \omterm{def:b3:topology:topology}{topologia} da norma em $\ell^1$ (as \omterm{def:b3:topology:topology}{topologias}, essas,
diferem --- \omterm{def:b3:topology:topology}{vizinhanças} fracas nunca são limitadas em norma), e
a \omterm{thm:b3:topology:metriccompact}{compacidade sequencial} fraca da bola unitária é uma propriedade diferente,
mais forte, equivalente à reflexividade
(Eberlein--Šmulian, além de nossas ferramentas; a falha, ao
menos, nós demonstramos).

\textbf{22.} \emph{O censo.}
\begin{center}
\begin{tabular}{l|l|c|c|c}
$E$ & $E'$ & sep. & refl. & bolas fraca.\ seq.\ cpt.\\
\hline
$c_0$ & $\ell^1$ & sim & não & não ($e_1{+}\dots{+}e_k$)\\
$\ell^1$ & $\ell^\infty$ & sim & não & não ($e_k$, q.~21)\\
$\ell^p$ & $\ell^q$ & sim & sim & sim (q.~13)\\
$\ell^\infty$ & $\supsetneq\ell^1$ & não & não & não\\
\end{tabular}
\end{center}
Assinaturas: $c_0$ --- seu bidual é $\ell^\infty$: o primeiro
passo não \omterm{rem:b3:banach:bidual}{reflexivo} (questão 14); $\ell^1$ --- a propriedade de
Schur (questão 20); $\ell^p$ --- reflexividade e \omterm{def:b3:topology:compact}{compacidade}
fraca (questões 7 e 13); $\ell^\infty$ ---
não separabilidade e limites de Banach: funcionais que sequência alguma
representa (questões 8 e 10--11). Uma família de espaços,
quatro mundos diferentes.

\textbf{23.} Seja $(f_k) \subseteq F$ com $\norm{x - f_k}
\to d$. Então $\norm{f_k} \leq \norm x + \sup_k\norm{x -
f_k}$: limitada, de modo que, pela questão 13, uma subsequência $f_{k_j}
\rightharpoonup f$. Se $f \notin F$, então $\delta =
\operatorname{dist}(f, F) > 0$ ($F$ fechado); em $F \oplus \R
f$, a forma linear $\lambda(g + tf) = t$ satisfaz
$\abs{\lambda(u)} \leq \norm u/\delta$ (pois $\norm{g +
tf} \geq \abs t\,\delta$), e Hahn--Banach a estende a
$\Lambda \in (\ell^p)'$ com $\Lambda\restriction_F = 0$,
$\Lambda(f) = 1$; mas então $0 = \Lambda(f_{k_j}) \to
\Lambda(f) = 1$: contradição. Logo $f \in F$, e $x -
f_{k_j} \rightharpoonup x - f$ dá, pela questão 17,
\[
d \leq \norm{x - f} \leq \liminf_j\,\norm{x - f_{k_j}} = d :
\]
a distância é atingida em $f$. Em $c_0$: $\abs{\Lambda(x)}
\leq \sum_n2^{-n}\abs{x_n} < \norm x_\infty$ para todo $x
\neq 0$ (uma sequência nula não nula não pode satisfazer $\abs{x_n} =
\norm x_\infty$ para todo $n$), ao passo que as truncadas $(1,
\dots, 1, 0, \dots)$ dão $\Lambda = 1 - 2^{-N} \to 1$: logo
$\norm\Lambda = 1$, nunca atingido. Fórmula da distância: para $f
\in \ker\Lambda$, $\abs{\Lambda(x)} = \abs{\Lambda(x - f)}
\leq \norm{x - f}$, de modo que $\operatorname{dist}(x, \ker\Lambda)
\geq \abs{\Lambda(x)}$; reciprocamente, para $u$ na bola unitária
com $\Lambda(u) \geq 1 - \varepsilon$, o vetor $f = x -
\frac{\Lambda(x)}{\Lambda(u)}u$ está em $\ker\Lambda$ com
$\norm{x - f} \leq \frac{\abs{\Lambda(x)}}{1 - \varepsilon}$:
igualdade. Se algum $f \in \ker\Lambda$ a atingisse, $z = x -
f$ satisfaria $\abs{\Lambda(z)} = \abs{\Lambda(x)} =
\norm z \neq 0$, de modo que $\Lambda$ atingiria sua norma em
$z/\norm z$: impossível. Um hiperplano fechado de $c_0$ sem
pontos mais próximos em lugar algum --- a reflexividade não era decorativa.

\textbf{24.} Ponha $y_1 = x^{(1)}$. Dado $y_1, \dots, y_j$,
cada aplicação $k \mapsto \langle y_i, x^{(k)}\rangle$ tende a $0$
($y_i \in \ell^2 = (\ell^2)'$), de modo que existe $k_{j+1}$ além
do índice anterior com $\abs{\langle y_i,
x^{(k_{j+1})}\rangle} \leq \frac1{j+1}$ para $i = 1, \dots,
j$; chame a escolha de $y_{j+1}$. Então
\[
\Bigl\lVert\sum_{j=1}^my_j\Bigr\rVert_2^2
= \sum_{j=1}^m\norm{y_j}_2^2
+ 2\sum_{j=2}^m\sum_{i<j}\langle y_i, y_j\rangle
\leq mC^2 + 2\sum_{j=2}^m\frac{j-1}j
\leq mC^2 + 2m,
\]
e, dividindo por $m^2$: $\norm{\frac1m\sum_jy_j}_2^2 \leq
\frac{C^2 + 2}m \to 0$. (Para um limite fraco $x \neq 0$, aplique
isso a $x^{(k)} - x$.) Na base ortonormal $(e_k)$ nem sequer é preciso
extrair: $\norm{\frac1m(e_1 + \dots +
e_m)}_2 = \frac{\sqrt m}m = \frac1{\sqrt m}$. As médias
convertem convergência fraca em convergência em norma: a
propriedade de Banach--Saks de $\ell^2$.

\textbf{25.} $A_mx = \frac1m\sum_{i=0}^{m-1}S^ix$, de modo que
a linearidade e a invariância por deslocamento dão $L(A_mx) = L(x)$. Para qualquer
$u$ limitada e $\varepsilon > 0$, tome $N$ com $u_n \leq
\limsup u + \varepsilon$ para $n \geq N$; a positividade aplicada
a $(\limsup u + \varepsilon)\mathbf 1 - S^Nu \geq 0$ e
$L(S^Nu) = L(u)$ dá $L(u) \leq \limsup u + \varepsilon$
e, simetricamente, $L(u) \geq \liminf u - \varepsilon$: logo
$\liminf A_mx \leq L(x) \leq \limsup A_mx$ para todo $m$. Se
$x$ é $T$-periódica, $A_Tx$ é a sequência constante igual à
média do período $\mu$: $L(x) = \mu$ para \emph{todo} limite
de Banach --- $\frac13$ em $(1,0,0,\dots)$, $\frac12$ em
$(0,1,0,1,\dots)$, como na questão 12. Para a sequência de blocos:
em $N = 3^j$ com $j$ par, o último bloco é só de uns, de modo que a
média de Cesàro vale $\geq \frac{3^j - 3^{j-1}}{3^j} = \frac23$;
em $N = 3^j$ com $j$ ímpar, todos os uns estão em $\intoc0
{3^{j-1}}$, de modo que a média vale $\leq \frac13$. Logo $p(x) \geq
\frac23$ e $-p(-x) = \liminf_n\frac{x_1 + \dots + x_n}n
\leq \frac13$. Em $M = c \oplus \R x$, defina $\Lambda_+(y +
tx) = \lim y + t\,p(x)$. Dominação por $p$: para $t > 0$,
a sublinearidade dá $p(tx) \leq p(y + tx) + p(-y)$, isto é,
$p(y + tx) \geq t\,p(x) + \lim y$ (note $p(\pm y) = \pm\lim
y$ para $y \in c$: as médias de Cesàro de uma convergente
convergem para seu limite); para $t = -s < 0$, $p(y) \leq p(y -
sx) + p(sx)$ dá $p(y - sx) \geq \lim y - s\,p(x)$; para $t
= 0$ há igualdade. Logo $\Lambda_+ \leq p$ em $M$, e
Hahn--Banach o estende a $L_+ \leq p$ em $\ell^\infty_\R$,
que é um limite de Banach exatamente como na questão 10 (a dominação
por $p$ fornece a positividade, a invariância por deslocamento e o valor
$\lim$ em $c$), com $L_+(x) = p(x) \geq \frac23$. O mesmo
cálculo com $\Lambda_-(y + tx) = \lim y - t\,p(-x)$
(usando $p(-sx) \leq p(y - sx) + p(-y)$ para o caso $t = -s
< 0$) fornece um limite de Banach $L_-$ com $L_-(x) = -p(-x) \leq
\frac13$. Dois limites de Banach, uma sequência, dois valores:
fora do mundo periódico (e, mais geralmente, \emph{quase
convergente}), um limite de Banach é uma escolha genuína.
\end{solution}
